% Generated by tools/textbook/export_listings.js from tools/decks/mastery.js and tools/concepts/registry.json: do not edit.
\begin{chapstatus}{Stato di apprendimento di questo capitolo}
\masterychip{now}{DA COMPRENDERE OGGI}\enspace i multiplexer; le reti combinatorie; i codici; decoder ed encoder; comparatori e parità; l'ALU.\par\smallskip
\masterychip{use}{DA UTILIZZARE OGGI}\enspace il programma completo, gli strumenti di PRG32 (si usa, non si studia, in questo corso).\par\smallskip
\masterychip{later}{DA COMPRENDERE PIÙ AVANTI}\enspace macchine a stati finiti (Capitolo~\ref{ch:sequential}).
\end{chapstatus}
\begin{chapstatus}{Percorso del capitolo}
Alla fine del capitolo saprete:
\begin{enumerate}[nosep,leftmargin=2em]
\item distinguere una macchina combinatoria da una sequenziale
\item usare decodificatori, multiplexer, comparatori e parità
\item spiegare l'ALU come unità in parallelo e un multiplexer
\end{enumerate}
Si appoggia su: solo l'installazione del Capitolo~\ref{ch:setup}. Prossimo passo: Capitolo~\ref{ch:verm1p2}.
\par\smallskip\textbf{Guardando indietro, guardando avanti.} Nella Lezione 02 un'istruzione era un numero a 32 bit, e il processore in qualche modo sapeva che cosa farne. Oggi un decodificatore legge un campo di quel numero, funct3, e pilota il multiplexer che sceglie il risultato dell'ALU. La Lezione 19 descrive tutti i campi e il ciclo di prelievo, decodifica, esecuzione. \emph{Richiamo:} Quale parte della parola di un'istruzione dice all'ALU di sommare invece di calcolare un AND?
\end{chapstatus}
